
\addcontentsline{toc}{chapter}{Executive Summary}\vspace{-1cm}

%Border
\begin{tikzpicture}[remember picture, overlay]
	\draw[line width = 4pt] ($(current page.north west) + (0.75in,-0.75in)$) rectangle ($(current page.south east) + (-0.75in,0.75in)$);
\end{tikzpicture}

The objective of this project was to develop an intelligent Visual Inspection of Solar Farm system capable of automating the monitoring and maintenance process of large-scale solar energy installations. The project focused on integrating drone technology, computer vision, thermal imaging, and artificial intelligence to improve the efficiency, reliability, and accuracy of solar farm inspections. The system was designed to overcome the limitations of conventional manual inspection methods, which are often time-consuming, labor-intensive, and prone to human error.

The project involved the development of a complete inspection framework consisting of data acquisition, image processing, fault detection, thermal analysis, and monitoring modules. High-resolution RGB and infrared images were captured using drone-mounted cameras and processed using deep learning algorithms to identify faults such as panel cracks, dust accumulation, shading effects, hotspot formation, overheating components, and other operational anomalies. The system also incorporated a centralized dashboard for visualization, reporting, and maintenance decision support.

During the implementation phase, various technologies including Raspberry Pi, OpenCV, Python, thermal imaging techniques, and YOLO-based deep learning models were utilized. Experimental testing was carried out using visual and thermal datasets under different operating conditions to evaluate fault detection performance and inspection efficiency. The developed system successfully demonstrated automated fault identification, thermal anomaly detection, and rapid report generation for maintenance planning.

The project resulted in a scalable and cost-effective inspection solution capable of reducing inspection time, minimizing manual effort, and improving the operational reliability of solar farms. In addition to technical knowledge in artificial intelligence, image processing, and renewable energy monitoring, the project contributed significantly to the development of analytical thinking, problem-solving abilities, teamwork, and project management skills. This report documents the design methodology, implementation process, experimental evaluation, and key outcomes achieved during the development of the Visual Inspection of Solar Farm system.

\pagebreak

